\documentclass[12pt,a4paper]{article}

\usepackage[margin=2cm]{geometry}
\usepackage{amsmath}
\usepackage{enumitem}

\setlength{\parindent}{0pt}
\setlength{\parskip}{6pt}

\begin{document}

\begin{center}
    {\Large\textbf{Statics -- Resolving Forces}}\\[4pt]
    \textbf{A-Level Physics}
\end{center}

For each question, draw a force diagram and resolve the forces
horizontally and vertically. Assume all objects are in equilibrium.

Where required, take \(g = 9.81\,\mathrm{m\,s^{-2}}\).

\vspace{0.5cm}

\begin{enumerate}[leftmargin=*, itemsep=1em]

\item
A sign of weight \(180\,\mathrm{N}\) is suspended by two cables.
One cable makes an angle of \(35^\circ\) above the horizontal,
and the other makes an angle of \(50^\circ\) above the horizontal.

Calculate the tension in each cable.

\vspace{1.2cm}

\item
A small ring is in equilibrium under three forces.
One force is \(80\,\mathrm{N}\) horizontally to the right.
A second force is \(120\,\mathrm{N}\) at \(65^\circ\) above
the horizontal to the left.

Calculate the \textbf{magnitude and direction} of the third force.

\vspace{1.2cm}

\item
A \(25\,\mathrm{kg}\) stage light is supported by two cables.
The left cable makes an angle of \(40^\circ\) above the horizontal,
while the right cable makes an angle of \(25^\circ\) above the horizontal.

Calculate the tension in each cable.

\vspace{1.2cm}

\item
A \(30\,\mathrm{kg}\) box rests on a horizontal floor.
A force of \(100\,\mathrm{N}\) is applied at \(30^\circ\)
below the horizontal. The box remains stationary.

Calculate:

\begin{enumerate}[label=(\alph*)]
    \item the normal reaction force;
    \item the frictional force.
\end{enumerate}

\vspace{1.2cm}

\item
A ring is acted upon by three forces.
A \(90\,\mathrm{N}\) force acts horizontally to the left and
a \(140\,\mathrm{N}\) force acts at \(55^\circ\) above the
horizontal to the right.

The third force keeps the ring in equilibrium.

Calculate the \textbf{magnitude and direction} of the third force.

\vspace{1.2cm}

\item
A \(300\,\mathrm{N}\) object is held stationary by two ropes.
One rope acts at \(20^\circ\) above the horizontal to the left,
and the other at \(60^\circ\) above the horizontal to the right.

Calculate the tension in each rope.

\end{enumerate}

\newpage

\begin{center}
    {\Large\textbf{Statics -- Answers}}
\end{center}

\vspace{0.5cm}

\begin{enumerate}[leftmargin=*, itemsep=1em]

\item
\(35^\circ\) cable: \(116\,\mathrm{N}\) \\
\(50^\circ\) cable: \(148\,\mathrm{N}\)

\item
Third force magnitude: \(113\,\mathrm{N}\) \\
Direction: \(74.9^\circ\) below the horizontal to the left

\item
Left cable tension: \(245\,\mathrm{N}\) \\
Right cable tension: \(207\,\mathrm{N}\)

\item
Normal reaction: \(344\,\mathrm{N}\) \\
Frictional force: \(86.6\,\mathrm{N}\)

\item
Third force magnitude: \(115\,\mathrm{N}\) \\
Direction: \(85.2^\circ\) below the horizontal to the right

\item
\(20^\circ\) rope tension: \(152\,\mathrm{N}\) \\
\(60^\circ\) rope tension: \(286\,\mathrm{N}\)

\end{enumerate}

\end{document}